%%% FIELD proof-prop-polynomial-phase-diagram
By the polynomial-radius specialization in the statement and the symbol declarations, \(n\in\mathbb N\) and \(a_0,a_1,a_e\geq 0\).  Hence
\[
0<r_{jn}=n^{-a_j}\leq 1\leq \bar r,
\qquad j\in\{0,1,e\},
\]
so the three radius bounds required by the bounded-radius regime hold.

We first extract the polynomial order of every summand in \(\rho_n\).  Define, only within this proof,
\[
t_\kappa:=\frac{\kappa}{2}+(1-\kappa)a_1,
\qquad
u_\kappa:=a_1+\frac{\kappa}{2(\kappa+1)},
\qquad
w_\kappa:=a_1+\frac{\kappa a_e}{\kappa+2},
\qquad
c:=a_0+a_e.
\]
The definition of \(Q_\kappa\) and the equality \(r_{1n}=n^{-a_1}\) give, when \(0<\kappa<1\),
\[
Q_\kappa(n,r_{1n})
=n^{-1/2}\vee n^{-\kappa/2}(n^{-a_1})^{1-\kappa}
=n^{-1/2}\vee n^{-t_\kappa}.
\]
The remaining three terms in the definition of \(\rho_n\) satisfy
\[
r_{1n}n^{-\kappa/[2(\kappa+1)]}=n^{-u_\kappa},
\qquad
r_{1n}r_{en}^{\kappa/(\kappa+2)}=n^{-w_\kappa},
\qquad
r_{0n}r_{en}=n^{-c}.
\]
For \(n\geq1\), the largest of finitely many powers \(n^{-s}\) is the power having the smallest exponent.  Therefore Definition \(\mathrm{def{:}phase\mbox{-}exponent}\) gives, for \(0<\kappa<1\),
\[
\beta_\kappa=\min\{1/2,t_\kappa,u_\kappa,w_\kappa,c\}.
\]
For \(\kappa>1\), Definition \(\mathrm{def{:}tail\mbox{-}noise\mbox{-}rate}\) gives \(Q_\kappa(n,r_{1n})=n^{-1/2}\), and hence
\[
\beta_\kappa=\min\{1/2,u_\kappa,w_\kappa,c\}.
\]
In either case, the summand whose exponent equals \(\beta_\kappa\) supplies the lower bound, while each of the four summands in Definition \(\mathrm{def{:}sharp\mbox{-}rate}\) is at most \(n^{-\beta_\kappa}\).  Thus
\[
n^{-\beta_\kappa}\leq \rho_n\leq 4n^{-\beta_\kappa},
\qquad \kappa\neq1.
\]

At \(\kappa=1\), put
\[
L_n:=\sqrt{1+\log_+(nr_{1n}^2)}.
\]
Then \(Q_1(n,r_{1n})=n^{-1/2}L_n\), \(L_n\geq1\), and
\[
\beta_1=\min\{1/2,a_1+1/4,a_1+a_e/3,c\}.
\]
The same comparison, now using \(L_n\geq1\), yields
\[
n^{-\beta_1}\leq \rho_n\leq4L_n n^{-\beta_1}.
\]
This proves the asserted polynomial order, including the stated possible boundary logarithm.

It remains to characterize strict deterioration relative to \(B\).  First suppose \(B\leq a_1\).  For \(0<\kappa<1\),
\[
t_\kappa=(1-\kappa)a_1+\kappa/2\geq B,
\]
because both \(a_1\) and \(1/2\) are at least \(B\).  For every \(\kappa>0\),
\[
u_\kappa\geq a_1\geq B,
\qquad
w_\kappa\geq a_1\geq B.
\]
Also \(1/2\geq B\) and \(c=a_0+a_e\geq B\) by the definition
\(B=\min\{1/2,a_1+a_e,a_0+a_e\}\).  Thus every entry defining \(\beta_\kappa\) is at least \(B\).  Conversely, if the minimum defining \(B\) is attained by \(1/2\) or by \(c\), that same entry occurs in \(\beta_\kappa\).  In the remaining case \(B=a_1+a_e\); the inequalities \(B\leq a_1\) and \(a_e\geq0\) then force \(a_e=0\) and \(B=a_1=w_\kappa\).  Hence
\[
\beta_\kappa=B
\]
whenever \(B\leq a_1\), proving that there is no polynomial deterioration in this case.

Now suppose \(B>a_1\), and set \(x:=B-a_1>0\) as in Definition \(\mathrm{def{:}deterioration\mbox{-}threshold}\).  Since \(B\leq1/2\) and \(B\leq a_1+a_e\),
\[
0<x\leq 1/2-a_1,
\qquad
0<x\leq a_e.
\]
In particular, \(a_1<1/2\), the first denominator below is positive, and \(a_e-x\geq0\).  Define the three extended-real cutoffs
\[
A_1:=\frac{x}{1/2-a_1},
\qquad
A_2:=\frac{2x}{1-2x},
\qquad
A_3:=\frac{2x}{a_e-x},
\]
with a zero denominator interpreted as \(+\infty\), exactly as in Definition \(\mathrm{def{:}deterioration\mbox{-}threshold}\).  Notice that \(A_1\leq1\).

The tail-noise exponent is present only for \(0<\kappa<1\), and direct subtraction gives
\[
t_\kappa<B
\quad\Longleftrightarrow\quad
\kappa(1/2-a_1)<x
\quad\Longleftrightarrow\quad
\kappa<A_1.
\]
Because \(A_1\leq1\), the condition \(0<\kappa<A_1\) automatically lies in the range where this entry is present.  For the empty-shell exponent, if \(x<1/2\), multiplication by the positive number \(2(\kappa+1)\) gives
\[
u_\kappa<B
\quad\Longleftrightarrow\quad
\frac{\kappa}{2(\kappa+1)}<x
\quad\Longleftrightarrow\quad
\kappa(1-2x)<2x
\quad\Longleftrightarrow\quad
\kappa<A_2.
\]
If \(x=1/2\), the strict inequality \(\kappa/[2(\kappa+1)]<x\) holds for every finite \(\kappa>0\), and the convention \(A_2=+\infty\) gives the same equivalence.  Finally, if \(a_e>x\),
\[
w_\kappa<B
\quad\Longleftrightarrow\quad
\frac{\kappa a_e}{\kappa+2}<x
\quad\Longleftrightarrow\quad
\kappa(a_e-x)<2x
\quad\Longleftrightarrow\quad
\kappa<A_3.
\]
If \(a_e=x\), then \(\kappa a_e/(\kappa+2)<a_e=x\) for every finite \(\kappa>0\), again agreeing with \(A_3=+\infty\).

The two entries \(1/2\) and \(c\) are at least \(B\) by definition.  Consequently, the preceding three displays exhaust all ways in which an entry of \(\beta_\kappa\) can be strictly smaller than \(B\).  Taking their union,
\[
\begin{aligned}
\beta_\kappa<B
&\quad\Longleftrightarrow\quad
(\kappa<A_1)\ \text{or}\ (\kappa<A_2)\ \text{or}\ (\kappa<A_3)\\
&\quad\Longleftrightarrow\quad
0<\kappa<\max\{A_1,A_2,A_3\}
\quad\Longleftrightarrow\quad
0<\kappa<\kappa_\dagger,
\end{aligned}
\]
where the last equivalence is Definition \(\mathrm{def{:}deterioration\mbox{-}threshold}\).  The strict inequalities also show that a finite boundary \(\kappa=\kappa_\dagger\) is not part of the deterioration region.

For completeness, the logarithmic qualification at \(\kappa=1\) is exact.  Since \(r_{1n}=n^{-a_1}\),
\[
L_n=\sqrt{1+\log_+(n^{1-2a_1})}.
\]
If \(a_1\geq1/2\), then \(n^{1-2a_1}\leq1\) and \(L_n=1\).  If \(a_1<1/2\), then \(L_n\) is of order \(\sqrt{\log n}\).  When \(\beta_1<1/2\), writing \(d:=1/2-\beta_1>0\) shows
\[
\frac{n^{-1/2}L_n}{n^{-\beta_1}}=n^{-d}L_n\longrightarrow0;
\]
indeed, after setting \(s=\log n\), the bound \(e^{ds}\geq (ds)^2/2\) proves the limit.  Hence the logarithm does not change the leading rate in that case.  When \(\beta_1=1/2\), all three non-noise terms are at most \(n^{-1/2}\), so
\[
n^{-1/2}L_n\leq\rho_n\leq4n^{-1/2}L_n.
\]
Thus the boundary factor changes a leading root-\(n\) rate exactly when \(a_1<1/2\) and \(\beta_1=1/2\), which in particular proves the necessary condition stated in the proposition.

Finally,
\[
\frac{d}{d\kappa}u_\kappa=\frac{1}{2(\kappa+1)^2}>0,
\qquad
\lim_{\kappa\to\infty}u_\kappa=a_1+1/2\geq1/2\geq B.
\]
Therefore the empty-shell channel has no surviving deterioration in the fast-tail limit; any deterioration at finite \(\kappa\) is already exactly accounted for by the cutoff \(A_2\) above.
